\section{Related Work}
\label{sec:related}
\subsection{Strictly Paired IVIF}
Most IVIF approaches are trained using strictly paired datasets. 
Depending on the deep learning framework employed, these methods can broadly be categorised into three types: CNN framework \cite{densefuse, rfn-nest, mufusion, mmdrfuse, lrrnet}, CNN-Transformer framework \cite{swinfusion, cddfuse, emma, AFTfusion}, and generative framework \cite{fusiongan, tardal, tgfuse, ddfm}.
To facilitate end-to-end training, \cite{densefuse}, \cite{rfn-nest} introduce dense blocks and residual blocks to explore multi-scale features and avoid the design of hand-crafted fusion rules. 
\cite{mmdrfuse} continues this end-to-end manner of training with a mini-model involving only two CNN blocks, demonstrating that complex network architectures are not necessary for pursuing high performance.
To emphasise the global long-range context, \cite{swinfusion} conducts fusion based on a cross-attention mechanism.
% introduces intra-domain and inter-domain fusion based on a cross-attention mechanism.
\cite{cddfuse} and \cite{emma} adopt an improved transformer block\cite{restormer} to achieve a trade-off between efficiency and performance.
% \rr{To better discriminate the modality-specific features, \cite{tardal, tgfuse} introduce a generative adversarial game with the generator producing a fused image and two discriminators gauging its quality from the perspective of each modality, helping the training of the generator.}    
To discriminate the modality-specific features, \cite{tardal, tgfuse} introduce two discriminators that help to enhance the modality-related clues.
However, GAN-based approaches confront issues such as mode collapse and training instability.
Although \cite{ddfm} applies a likelihood rectification that achieves stable training without a mode collapse, it is computationally heavy.  

In this work, to cover a wide scope in terms of IVIF methodologies, we construct three lightweight baseline models to verify the feasibility, degree of unpairedness, and the merits of APTP and UPTP. 
% To validate the universality of this mode of training, we build and evaluate stand-alone  CNN and Transformer frameworks.
\subsection{Unpaired Training Mode}
% \rr{Unpaired training mode is first introduced in the image-to-image translation task to address the prohibitive expense of collecting large-scale paired training datasets \cite{pix2pix}. 
% In particular, \cite{cyclegan} breaks this limitation by introducing adversarial and cycle-consistent supervisions, achieving style transfer with consistent content.  
% This innovative paradigm enables training with unpaired data, which greatly broadens task applications. 
% Such challenging problems also exist in the field of image restoration and IVIF. 
% Following \cite{cyclegan}, \cite{unpaire_restoration1}, \cite{unpaire_restoration2} frame restoration as an unpaired style-transfer task from degraded to clean images. 
% Unlike translation or restoration, IVIF demands consistency between two complementary images rather than within a single one, and lacks a well-defined target domain, rendering unpaired training like \cite{cyclegan} infeasible. 
% Yet it is precisely the absence of a well-defined target domain, coupled with the pixel-level nature of the IVIF task, that makes unpaired training still feasible. }
% % Motivated by the independent analysis in \cite{unpaire_restoration3}, 
Unpaired training is first introduced in the image generation field to address the prohibitive expense of collecting large-scale paired datasets \cite{pix2pix}. 
\cite{cyclegan} breaks this limitation by introducing adversarial and cycle-consistent supervisions, achieving style transfer with consistent content.  
This innovative paradigm enables training with unpaired data, which greatly broadens task applications. 
Following \cite{cyclegan}, \cite{unpaire_restoration1}, \cite{unpaire_restoration2} frame restoration as an unpaired generation task from degraded to clean images. 
Such challenging problems also exist in the field of IVIF.
Unlike image generation, IVIF demands content consistency between two complementary images rather than within a single one, and lacks a well-defined fused image domain, rendering unpaired training like \cite{cyclegan} infeasible. 
Yet it is precisely the absence of a desired fused image, coupled with the pixel-level nature of the IVIF task, that makes unpaired training still feasible.


Due to the absence of ground-truth, fusion networks are not required to learn from an ideal fused image. 
Instead, supervision is usually drawn from the source images.
Therefore, we argue that the model can still extract the corresponding supervisory signals, even if the input is unpaired. 
% Furthermore, since the IVIF task is performed at the pixel level, the model can still learn to fuse multi-modal images by integrating cross-modal pixels, even when the source-modal pixels do not correspond to the same object. 
Starting from these points, we pioneer in exploring the possibility of training with unpaired source images. 
